
%!TEX root = main.tex
\usepackage{shellesc}

\usepackage{amsmath}
\usepackage{url}
\usepackage{graphicx}
\usepackage{siunitx}									
\usepackage{lipsum}
\usepackage{enumitem}
\usepackage{xcolor}
\usepackage{xfrac}	
% \usepackage{amssymb}			
\usepackage{upgreek}				
\usepackage{icomma}	  						
\usepackage{trsym}											
\usepackage{trfsigns}
%\usepackage{subcaption}
\usepackage[right]{eurosym}					
\usepackage{textcomp}									
			
			
%\usepackage[sorting=none, url=false]{biblatex}
			
\newcommand{\tsup}{\textsuperscript}
\newcommand{\tsub}{\textsubscript}
\newcommand{\mi}[2]{\ensuremath{#1_{\text{#2}}}}
\newcommand{\mis}[3]{\ensuremath{#1_{\text{#2}}^{\text{#3}}}}
\newcommand{\mtis}[3]{\ensuremath{\text{#1}_{\text{#2}}^{\text{#3}}}}
\newcommand{\mti}[2]{\ensuremath{\text{#1}_{\text{#2}}}}
\newcommand{\ddt}[1]{\ensuremath{\frac{\textnormal{d}#1}{\textnormal{d}t}}}
\newcommand{\sddt}[1]{\ensuremath{\sfrac{\textnormal{d}#1}{\textnormal{d}t}}}
\newcommand{\fref}[1]{Fig.~\ref{#1}}


% \newcommand{\linebreakand}{%
%     \end{@IEEEauthorhalign}
%     \hfill\mbox{}\par
%     \mbox{}\hfill\begin{@IEEEauthorhalign}
% }


\usepackage{tikz}
\usetikzlibrary{external}
\tikzexternalize[prefix=tikz_extern/]

%\usepackage{subfig}
%\usepackage{subcaption}
%\usepackage{subfloat}

 % ,pgf,pgfplots}					%% Auch wenn tikz kein Zeichenprogramm ist, es macht schöne Grafiken
%\usepgfplotslibrary{groupplots}

\graphicspath{{./figures/}}
\makeatletter
\def\input@path{{./figures/}}

\usepackage{pgf,pgfplots}					%% Auch wenn tikz kein Zeichenprogramm ist, es macht schöne Grafiken
%\usepgfplotslibrary{groupplots} %\usepgfplotslibrary{external} 
%\usepackage[european,straightvoltages,siunitx]{circuitikz}
\pgfplotsset{compat=newest}						%% behebt einige Fehler, die beim Rendern der Grafiken auftreten können
\usetikzlibrary{plotmarks,shapes,positioning,fit}
%\pgfkeys{/pgf/number format/.cd ,use comma ,set thousands separator={.}} %Globale einstellungen für Grafiken mit tikz/pgf.
%\usetikzlibrary{spy} 							%% kann bei Bedarf eingebunden werden, um Lupeneffekt bei Grafiken zu ermöglichen
\tikzset{every picture/.style={line join=round}}

%Doku matlab2tikz:
\pgfplotsset{plot coordinates/math parser=false} 

%sets search directory for pgfplots input files 
%\pgfplotsset{table/search path={tikz/tikzData}}
%
%\pgfplotsset{
%	every axis plot post/.style={
%		line join=round
%	}
%}


%% Parameter für Grafikgröße (tikz)
\newlength\figureheight
\newlength\figurewidth 

%% Folgende Zeile einkommentieren, um externalize zu aktivieren und alle mit tikz erstellten Grafiken als pdf auszugeben.
%% -> Zwischenspeicherung der gerenderten Bilder ermöglicht schnelleres kompilieren
%\usepgfplotslibrary{external}						% this one kills it with lulatex
%\pgfkeys{/pgf/images/include external/.code=\includegraphics{#1}}

\usepackage{filemod}

%\ifdef{\tikzexternalize}{
%\tikzexternalize[prefix=imgs/]					%% Zwischenspeicher im Unterordner imgs/ des tex/-Ausgabeordners
%\newcommand{\includetikz}[1]
%{
%	\tikzsetnextfilename{#1}
%	\filemodCmp{tikz/#1.tex}{imgs/#1.pdf}
%	{\tikzset{external/remake next}}{}
%	\input{tikz/#1.tex}
%}
%}
%{
%	\newcommand{\includetikz}[1]{
%		\input{tikz/#1.tex}
%	}
%}
%
%\input{kob_pgfplots}
%\input{rwth-colors}

\definecolor{mycolor}{RGB}{33, 95, 139}
            
\definecolor{BoschRGB1} {rgb}{0.8863    ,0    	  ,0.0824 }
\definecolor{BoschRGB2} {rgb}{0.7255    ,0.0078    ,0.4627}
\definecolor{BoschRGB3} {rgb}{0.3137    ,0.1373    ,0.4980}
\definecolor{BoschRGB4} {rgb}{0    		,0.3373    ,0.5686}
\definecolor{BoschRGB5} {rgb}{0    		,0.5569    ,0.8118}
\definecolor{BoschRGB6} {rgb}{0    		,0.6588    ,0.6902}
\definecolor{BoschRGB7} {rgb}{0.4706    ,0.7451    ,0.1255}
\definecolor{BoschRGB8} {rgb}{0   		,0.3843    ,0.2863}
\definecolor{BoschRGB9} {rgb}{0.3216    ,0.3725    ,0.4196}
\definecolor{BoschRGB10}{rgb}{0.7529    ,0.7529    ,0.7608}  
\definecolor{BoschRGB11}{rgb}{0     	,0     	  ,0}  
\definecolor{BoschRGB12}{rgb}{1			,1 		  ,1}